\section*{Chapitre \ref{ch:g4:volcanoes-earthquakes-first-look} --- Volcans et séismes~: un premier regard}

\begin{solution}{exo:g4:volcanoes-earthquakes-first-look:1}
(a) De la \omterm{def:g4:volcanoes-earthquakes-first-look:eruption}{lave}. (b) Des \omterm{def:g4:volcanoes-earthquakes-first-look:eruption}{cendres volcaniques}. (c) Des gaz.
\end{solution}

\begin{solution}{exo:g4:volcanoes-earthquakes-first-look:2}
Un \omterm{def:g4:volcanoes-earthquakes-first-look:earthquake}{séisme} est une secousse brutale du sol, causée par des \omterm{def:g3:stones-pebbles-sand:rock}{roches} qui se
cassent et glissent en profondeur. Il dure quelques secondes (une minute ou
deux pour les plus forts), pas des heures.
\end{solution}

\begin{solution}{exo:g4:volcanoes-earthquakes-first-look:3}
Se baisser, s'abriter, s'agripper.
\end{solution}

\begin{solution}{exo:g4:volcanoes-earthquakes-first-look:4}
Autour de l'\omterm{def:g2:water-on-earth:ocean}{océan} Pacifique.
\end{solution}

\begin{solution}{exo:g4:volcanoes-earthquakes-first-look:5}
(a) Vrai~: seuls les instruments remarquent la plupart d'entre eux.
(b) Faux~: la \omterm{def:g4:volcanoes-earthquakes-first-look:eruption}{lave} se refroidit et durcit en une \omterm{def:g3:stones-pebbles-sand:rock}{roche} nouvelle. (c) Faux~:
courir pendant que le sol tremble est dangereux (on peut tomber, et des
éclats de verre et des objets tombent des bâtiments)~; il faut se baisser,
s'abriter et s'agripper.
\end{solution}

\begin{solution}{exo:g4:volcanoes-earthquakes-first-look:6}
Elle vient des profondeurs, sous le \omterm{def:g4:volcanoes-earthquakes-first-look:volcano}{volcan}, là où la \omterm{def:g3:stones-pebbles-sand:rock}{roche} a fondu. Elle
monte par le milieu du \omterm{def:g4:volcanoes-earthquakes-first-look:volcano}{volcan} et sort au sommet.
\end{solution}

\begin{solution}{exo:g4:volcanoes-earthquakes-first-look:7}
(a) Le Chili (le Pérou, l'Équateur ou le Mexique conviennent aussi)~: ses
\omterm{def:g4:volcanoes-earthquakes-first-look:volcano}{volcans} s'alignent sur toute sa longueur, près de la \omterm{def:g1:landscapes-around-us:water}{côte} du Pacifique.
(b) Hawaï, au milieu de l'\omterm{def:g2:water-on-earth:ocean}{océan} Pacifique, loin de l'anneau ombré.
\end{solution}

\begin{solution}{exo:g4:volcanoes-earthquakes-first-look:8}
$141 \div 12$ fait un peu moins de 12 (car $12 \times 12 = 144$)~: environ
12 grands \omterm{def:g4:volcanoes-earthquakes-first-look:earthquake}{séismes} par mois.
\end{solution}

\begin{solution}{exo:g4:volcanoes-earthquakes-first-look:9}
Sur 20, un quart fait 5, donc trois quarts font $3 \times 5 = 15$~:
environ 15 dans l'anneau, et environ $20 - 15 = 5$ en dehors.
\end{solution}

\begin{solution}{exo:g4:volcanoes-earthquakes-first-look:10}
Il aurait dû se baisser sur les mains et les genoux, protéger sa tête et sa
nuque sous sa table, et s'y agripper jusqu'à la fin des secousses.
L'encadrement d'une porte ne le protège pas des objets qui tombent ou qui
sont projetés, et courir pendant que le sol tremble peut le faire tomber.
\end{solution}

\begin{solution}{exo:g4:volcanoes-earthquakes-first-look:11}
En dehors de l'anneau~: $1214 - 837 = 377$ \omterm{def:g4:volcanoes-earthquakes-first-look:volcano}{volcans}. Un tiers de 1\,200
fait 400, donc les deux tiers de 1\,200 font 800. L'anneau en contient 837,
un peu plus que 800 sur environ 1\,200~: «~environ 2 sur 3~» est une bonne
description.
\end{solution}
